\documentclass[10pt,a4paper]{amsart}
\usepackage[margin=19mm]{geometry}
\usepackage{amsmath,amssymb,mathtools}
\usepackage[scr=rsfs,cal=euler]{mathalfa}
\usepackage{xfrac}
\usepackage[unicode,hidelinks,bookmarks=false]{hyperref}
\newcommand{\ie}{i.e.\ }
\newcommand{\cf}{cf.\ }
\newcommand{\resp}{respectively\ }
\newcommand{\Subsection}{\subsection}
\providecommand{\nobreakdash}{\nobreak\hbox{-}}
\setlength{\emergencystretch}{2em}
\multlinegap0pt
\usepackage[shortlabels]{enumitem}
\usepackage{amssymb}
\usepackage{xcolor}
\newtheorem{proposition}{Proposition}
\newcommand{\UP}{\blacktriangle}
\title{Characterising regular double Stone and pseudocomplemented Kleene algebras in completions of rough sets induced by reflexive relations}
\author{Jouni J\"{a}rvinen, S\'{a}ndor Radeleczki}
\date{Source: arXiv:2603.26883v3 (CC BY 4.0)}
\begin{document}
\maketitle

\noindent\emph{Source and licence.} Characterising regular double Stone and pseudocomplemented Kleene algebras in completions of rough sets induced by reflexive relations. Version arXiv:2603.26883v3 (CC BY 4.0), distributed by arXiv under the Creative Commons Attribution 4.0 International licence, http://creativecommons.org/licenses/by/4.0/, as recorded in the arXiv licence metadata for this exact version and shown on the abstract page https://arxiv.org/abs/2603.26883v3. Full licence notice: \texttt{licenses/arxiv-2603.26883-CC-BY-4.0.txt}.\par\smallskip

\noindent\textit{Editorial scope.} Counted unit: the article's proposition prop:irredundant\_cover (source0.tex lines 1820--1827). The article's complete original proof of it is delivered with it (source0.tex lines 1830--1872, one \texttt{proof} environment of 43 lines). No mathematics is written, repaired or re-derived: the statement and the proof are the article's own lines, in the article's own notation, and no retained line is substituted or re-flowed.  Retained uncounted as the source-local context the proof needs: lines 1661--1665.  Nothing was omitted from any retained span, so the package carries no omission record.  No retained line contains a citation, so the wrapper carries no bibliography and no bibliography entry.  No retained line contains a figure environment, an image-inclusion command or drawing code, so no image is delivered and the package declares no asset dependency.  Editorial formatting only, in addition: an AMS \texttt{amsart} preamble that restates the article's own declarations and macros used by the retained lines, namely the environments \texttt{proposition} on separate counters, together with the standard packages those lines need.  The retained environments print in this document as Proposition 1 in order of appearance, whereas the article prints the same items with the numbering of its own full document.  The complete native source of the article as distributed in the arXiv e-print for this exact version is bundled unchanged as \texttt{sources/197307/source0.tex}.
\begin{equation} \tag{rSt} \label{eq:rSt}
\begin{minipage}[c]{0.9\textwidth}
 For any $x \in U$, $\{x\}^{\UP}$ is an atom in $\wp(U)^\UP$ and $\mathfrak{core}\breve{R}(x) \neq \emptyset$.
\end{minipage}
\end{equation}

\begin{proposition} \label{prop:irredundant_cover}
Let $R$ be a reflexive relation on $U$. Then the following are equivalent:
\begin{enumerate}[label={\rm (\roman*)}]
\item $\{ \{x\}^\UP \mid x \in U\}$ is an irredundant covering of $U$.

\item Condition \eqref{eq:rSt} holds.
\end{enumerate}
\end{proposition}

\begin{proof}
$\mathrm{(i)}\Rightarrow\mathrm{(ii)}$: Let $x \in U$.
As $\{ \{x\}^\UP \mid x \in U\}$ forms an irredundant 
covering of $U$, the inclusion 
$\{y\}^\UP \subseteq\{x\}^\UP$ implies 
$\{y\}^\UP = \{x\}^\UP$  for any $y \in U$.
Let $\emptyset \neq Y^\UP \subseteq \{x\}^\UP$. Then $Y$ is nonempty
and it contains at least one element $y$.
Now $\{y\}^\UP \subseteq Y^\UP \subseteq \{x\}^\UP$ and 
$\{y\}^\UP = \{x\}^\UP$ imply $Y^\UP = \{x\}^\UP$. This means that $\{x\}^\UP$ is an atom of $\wp(U)^\UP$.

Next we prove that 
$\mathfrak{core}\breve{R}(x) \neq \emptyset$.
Because the covering 
$\{ \{a\}^\UP \mid a \in U\}$ is irredundant,
there exists an element $y$ in the difference
\[
\{x\}^\UP \setminus   
\bigcup \{\{b\}^\UP \mid \{b\}^\UP \neq\{x\}^\UP\}.
\]
Suppose that $y \in \breve{R}(c) = \{c\}^\UP$
for some $c \in U$. We must have 
$\breve{R}(c) = \{c\}^\UP = \{x\}^\UP = \breve{R}(x)$.
Thus, $\breve{R}(x) \subseteq \breve{R}(c)$ yields 
$y \in \mathfrak{core}\breve{R}(x)$.


\medskip\noindent%
$\mathrm{(ii)}\Rightarrow\mathrm{(i)}$: 
Suppose that \eqref{eq:rSt} holds and 
that the covering $\{ \{a\}^\UP \mid a \in U\}$ is 
not irredundant. This means that there exists $x \in U$
such that 
$\{x\}^\UP \subseteq \bigcup\{\{y\}^\UP \mid \{y\}^\UP \neq\{x\}^\UP\}$. 
Now $\emptyset\neq\mathfrak{core}\breve{R}(x) \subseteq\breve{R}(x)
= \{x\}^\UP$. Hence, for any $w \in \mathfrak{core}\breve{R}(x)$,
there is $y\in U$ with $\{y\}^\UP \neq \{x\}^\UP$
and $w\in \{y\}^\UP = \breve{R}(y)$. Then $\{x\}^\UP = \breve{R}(x) 
\subseteq \breve{R}(y) = \{y\}^\UP$. Since
$\{y\}^\UP$ is an atom in $\wp(U)^\UP$, we get
$\{y\}^\UP = \{x\}^\UP$, a contradiction. Thus,
we have an irredundant covering.
\end{proof}

\end{document}
